\section{Economic Robustness and the Strengthened HJB Procedure}
\label{app:r16robustnessrecord}
This experiment changes the continuous-diffusion capital economy and trains
all procedures anew. It is separate from both the confirmation of 128
previously selected policies and the finite-period continuation menu. Its
four economic calibrations, dimensions ten and fifty, and eight fixed
training streams give 64 trials. The report retains all 256 construction
processes and 320 deployed policy outputs. All procedures completed without
reference-policy fallback. A fresh bank of 8,192 paths on 2,048 cells
assesses each output, with common innovations within a trial. The
648-event family has allocation $\alpha=0.01$ and includes method averages,
individual-policy endpoints, and the declared direct comparisons. No
calibration or stream is selected on its outcome.

\subsection{Economic variation}
Table~\ref{tab:r16ctcalibrations} gives the complete economic grid. The high
calibration retains the original high-volatility diffusion; low changes
its diffusion risk. The longer-horizon economy has horizon two years.
The joint-stress economy changes coupling, terminal dispersion costs,
initial heterogeneity, and the actor's feasible radius together. Its
comparison therefore describes that economic design rather than a
single-primitive causal effect. These names are local to the continuous
economy: in particular, the finite menu's Long calibration has a four-year
horizon and half-year periods.

\begin{table}[htbp]
\centering\small
\caption{Continuous-economy robustness designs}
\label{tab:r16ctcalibrations}
\begin{tabular}{lrrrrrr}
\toprule
Design & $T$ & $\sigma_I$ & $\sigma_C$ & $\kappa$ & $\chi$ & Actor radius\\
\midrule
high & 1 & .60 & .30 & .15 & .03 & .10\\
low & 1 & .15 & .10 & .15 & .03 & .10\\
long & 2 & .30 & .15 & .15 & .03 & .10\\
stress & 1 & .30 & .15 & .30 & .10 & .20\\
\bottomrule
\end{tabular}
\par\medskip\begin{minipage}{.97\linewidth}\footnotesize\raggedright
All designs use dimensions 10 and 50, eight fixed streams per dimension,
discount rate .04, productivity .10, adjustment parameter .20, and action
bounds $[.02,2]$. The first three designs use nine initial profiles from
means $\{-.5,0,.5\}$ and dispersions $\{0,.25,.5\}$. Stress uses
means $\{-.75,0,.75\}$ and dispersions $\{0,.375,.75\}$. The feasible
actor radius is measured around the reference schedule.
\end{minipage}
\end{table}

\subsection{HJB differentiation, selection, and deployment}
The two-layer tanh value network propagates its value, gradient, and
covariance-contracted Hessian by the chain rule. The diffusion trace is exact
for this network in both dimensions, so training uses no random trace
estimate. Independent automatic-differentiation checks compare both the
network derivatives and their parameter gradients. The source validation
records discrepancies at floating-point precision. Stochastic traces with
2, 8, and 32 probes are evaluated only on an untouched finite diagnostic
bank after selection, for both covariance-Gaussian probes and idiosyncratic
Rademacher probes with the common-noise term treated exactly.

Four candidates combine learning rates .0015 and .0005 with either
reference occupation collocation or an equal mixture of reference occupation
and broad collocation. Each uses up to 1,200 Adam updates and 150 L-BFGS
iterations. Checkpoints are evaluated on independent occupation and broad
residual banks; the combined residual criterion chooses the candidate and
checkpoint. Its prescribed residual threshold is .001. Final payoff
observations play no role in selection. Every candidate, including an
unsuccessful or unselected candidate, is charged to the procedure.

The selected value network yields two retained actions. The greedy action
uses the explicit first-derivative calculation and a bracketed scalar root
solver with Newton updates; the prescribed bisection fallback remains part
of the implementation. A width-32 actor is separately fitted to these
actions. Independent diagnostic states record its action difference and
Hamiltonian difference from the greedy action. These are finite-state
diagnostics. They do not establish the economic loss of distillation, since
the protected family does not contain a direct greedy-minus-distilled
payoff contrast.

The diagnostic tables below retain all four configurations in all eight
economic cells. Their means and ranges use all eight streams in each cell.
The complete checkpoint histories, time-stratified residuals, selected
network identities, root counters, and action diagnostics remain in the
unchanged machine-readable report. A zero terminal discrepancy follows
from the specified terminal lift. Residuals on finite independent banks and
changes between iterates are reported with their corresponding scope;
they are not uniform PDE residual bounds or errors against the unknown
optimal value and policy.

% Authored descriptive tables; all generating scientific records remain unchanged.
% Frozen REPORT SHA256: 2b675db41f2827d987c3ceae3e3ee01627a0e80228e0b9f54f85d2922d7e17fd
\begingroup\small
\setlength{\tabcolsep}{3pt}
\begin{longtable}{lrlrlrrr}
\caption{All prespecified HJB candidate configurations}\label{tab:r16hjbconfigurations}\\
\toprule
Design & $d$ & Config. & Mean RMS & RMS range & Chosen & Stop & Fit sec. \\
\midrule
\endfirsthead
\toprule
Design & $d$ & Config. & Mean RMS & RMS range & Chosen & Stop & Fit sec. \\
\midrule
\endhead
high & 10 & O15 & 0.006102 & [0.005151, 0.006967] & 0/8 & 0/8 & 14.10 \\
high & 10 & O05 & 0.006280 & [0.005199, 0.007189] & 0/8 & 0/8 & 14.08 \\
high & 10 & M15 & 0.005500 & [0.004842, 0.006109] & 6/8 & 0/8 & 14.10 \\
high & 10 & M05 & 0.005936 & [0.005037, 0.006863] & 2/8 & 0/8 & 14.19 \\
high & 50 & O15 & 0.004162 & [0.003410, 0.005158] & 1/8 & 0/8 & 35.96 \\
high & 50 & O05 & 0.006985 & [0.004926, 0.010107] & 0/8 & 0/8 & 35.75 \\
high & 50 & M15 & 0.003714 & [0.003368, 0.004140] & 7/8 & 0/8 & 34.98 \\
high & 50 & M05 & 0.004732 & [0.003850, 0.007696] & 0/8 & 0/8 & 34.61 \\
low & 10 & O15 & 0.002043 & [0.001504, 0.002651] & 2/8 & 0/8 & 15.88 \\
low & 10 & O05 & 0.002567 & [0.001879, 0.003504] & 0/8 & 0/8 & 15.90 \\
low & 10 & M15 & 0.001684 & [0.001263, 0.002372] & 4/8 & 0/8 & 15.61 \\
low & 10 & M05 & 0.002085 & [0.001273, 0.002628] & 2/8 & 0/8 & 15.75 \\
low & 50 & O15 & 0.002916 & [0.002202, 0.003680] & 3/8 & 0/8 & 37.97 \\
low & 50 & O05 & 0.004182 & [0.003633, 0.004568] & 0/8 & 0/8 & 38.28 \\
low & 50 & M15 & 0.002770 & [0.001780, 0.004655] & 5/8 & 0/8 & 36.42 \\
low & 50 & M05 & 0.004554 & [0.003358, 0.006472] & 0/8 & 0/8 & 37.06 \\
long & 10 & O15 & 0.006106 & [0.004142, 0.007295] & 0/8 & 0/8 & 15.73 \\
long & 10 & O05 & 0.006880 & [0.005568, 0.009664] & 0/8 & 0/8 & 15.97 \\
long & 10 & M15 & 0.004367 & [0.003966, 0.004945] & 8/8 & 0/8 & 15.50 \\
long & 10 & M05 & 0.005465 & [0.004689, 0.006388] & 0/8 & 0/8 & 15.82 \\
long & 50 & O15 & 0.004367 & [0.003918, 0.005369] & 0/8 & 0/8 & 35.13 \\
long & 50 & O05 & 0.009539 & [0.006777, 0.011023] & 0/8 & 0/8 & 35.27 \\
long & 50 & M15 & 0.003349 & [0.002680, 0.003739] & 8/8 & 0/8 & 32.72 \\
long & 50 & M05 & 0.010590 & [0.005321, 0.013372] & 0/8 & 0/8 & 33.10 \\
stress & 10 & O15 & 0.005248 & [0.004654, 0.006116] & 0/8 & 0/8 & 16.57 \\
stress & 10 & O05 & 0.006521 & [0.005367, 0.008260] & 0/8 & 0/8 & 16.55 \\
stress & 10 & M15 & 0.004643 & [0.003880, 0.005171] & 7/8 & 0/8 & 16.35 \\
stress & 10 & M05 & 0.005211 & [0.004528, 0.005862] & 1/8 & 0/8 & 16.22 \\
stress & 50 & O15 & 0.004790 & [0.004261, 0.005632] & 0/8 & 0/8 & 34.63 \\
stress & 50 & O05 & 0.006495 & [0.005779, 0.007554] & 0/8 & 0/8 & 34.49 \\
stress & 50 & M15 & 0.003732 & [0.003409, 0.004341] & 8/8 & 0/8 & 33.85 \\
stress & 50 & M05 & 0.005291 & [0.004535, 0.006117] & 0/8 & 0/8 & 33.73 \\
\bottomrule
\end{longtable}
\par\medskip\noindent\footnotesize Each row retains all eight declared streams. O15 and O05 use occupation training collocation with learning rates $0.0015$ and $0.0005$; M15 and M05 use equal occupation and broad training collocation at the same two rates. Selection always gives equal weight to the separate occupation and broad selection banks. Mean RMS and its observed range use the best retained checkpoint of each candidate. Chosen counts residual-only family winners; Stop counts residual-only early stopping. Fit seconds are per-candidate components, already included in the complete shared HJB work in Table~\ref{tab:r16-robustness-work}. All 256 candidate fits and 2048 saved checkpoints are retained; no candidate failed and no final payoff selected a checkpoint or configuration.
\endgroup
\begingroup\footnotesize
\setlength{\tabcolsep}{3pt}
\begin{longtable}{lrllll}
\caption{Untouched HJB residual audits and selected-checkpoint iterate changes}\label{tab:r16hjbaudit}\\
\toprule
Design & $d$ & Occupation RMS & Broad RMS & $\Delta V$ RMS & $\Delta u$ RMS \\
\midrule
\endfirsthead
\toprule
Design & $d$ & Occupation RMS & Broad RMS & $\Delta V$ RMS & $\Delta u$ RMS \\
\midrule
\endhead
high & 10 & [0.003482, 0.004507] & [0.005814, 0.007266] & [0.002822, 0.008041] & [0.006681, 0.014383] \\
high & 50 & [0.002561, 0.003223] & [0.003911, 0.005019] & [0.001268, 0.003759] & [0.007258, 0.019404] \\
low & 10 & [0.000969, 0.001512] & [0.001476, 0.002424] & [0.002016, 0.013412] & [0.012218, 0.032167] \\
low & 50 & [0.001409, 0.002877] & [0.002093, 0.004012] & [0.002262, 0.007162] & [0.029766, 0.046553] \\
long & 10 & [0.002812, 0.003629] & [0.005056, 0.006162] & [0.005533, 0.011884] & [0.009387, 0.019914] \\
long & 50 & [0.002297, 0.003004] & [0.003169, 0.004497] & [0.005593, 0.015276] & [0.020734, 0.040882] \\
stress & 10 & [0.002724, 0.003714] & [0.004856, 0.006165] & [0.002291, 0.007751] & [0.010873, 0.027689] \\
stress & 50 & [0.002594, 0.003421] & [0.003962, 0.005011] & [0.001553, 0.002530] & [0.014736, 0.022117] \\
\bottomrule
\end{longtable}
\par\medskip\noindent\footnotesize Each entry is the observed minimum and maximum over all eight streams in that economic cell. Residuals use the selected critic on untouched audit banks of 1024 occupation and 1024 broad states; these audit outcomes did not select a candidate. The $\Delta V$ and $\Delta u$ columns use the saved selection-bank RMS changes from Adam step 1200 to the selected final L-BFGS checkpoint, which was selected in all 64 trials. They measure differences between successive saved iterates. They do not measure error against the unknown optimal value or policy. Ranges are descriptive, not confidence intervals or uniform residual bounds.
\endgroup
\begingroup\footnotesize
\setlength{\tabcolsep}{3pt}
\begin{longtable}{lrlrrl}
\caption{Finite-bank trace-probe variance for both retained estimators}\label{tab:r16hjbprobes}\\
\toprule
Design & $d$ & Probe law & Probes & Mean variance & Variance range \\
\midrule
\endfirsthead
\toprule
Design & $d$ & Probe law & Probes & Mean variance & Variance range \\
\midrule
\endhead
high & 10 & Gaussian & 2 & 7.274e-05 & [6.700e-05, 7.862e-05] \\
high & 10 & Gaussian & 8 & 1.849e-05 & [1.671e-05, 1.972e-05] \\
high & 10 & Gaussian & 32 & 4.639e-06 & [4.152e-06, 5.020e-06] \\
high & 10 & Rademacher & 2 & 1.222e-05 & [1.049e-05, 1.492e-05] \\
high & 10 & Rademacher & 8 & 3.018e-06 & [2.602e-06, 3.691e-06] \\
high & 10 & Rademacher & 32 & 7.678e-07 & [6.502e-07, 9.196e-07] \\
high & 50 & Gaussian & 2 & 1.276e-05 & [1.099e-05, 1.494e-05] \\
high & 50 & Gaussian & 8 & 3.176e-06 & [2.696e-06, 3.804e-06] \\
high & 50 & Gaussian & 32 & 8.017e-07 & [6.780e-07, 9.383e-07] \\
high & 50 & Rademacher & 2 & 2.607e-06 & [1.509e-06, 3.877e-06] \\
high & 50 & Rademacher & 8 & 6.708e-07 & [3.858e-07, 9.657e-07] \\
high & 50 & Rademacher & 32 & 1.614e-07 & [9.272e-08, 2.294e-07] \\
low & 10 & Gaussian & 2 & 4.242e-07 & [3.127e-07, 6.243e-07] \\
low & 10 & Gaussian & 8 & 1.044e-07 & [7.834e-08, 1.548e-07] \\
low & 10 & Gaussian & 32 & 2.632e-08 & [1.934e-08, 3.989e-08] \\
low & 10 & Rademacher & 2 & 9.637e-08 & [4.980e-08, 1.793e-07] \\
low & 10 & Rademacher & 8 & 2.416e-08 & [1.230e-08, 4.672e-08] \\
low & 10 & Rademacher & 32 & 5.992e-09 & [3.043e-09, 1.190e-08] \\
low & 50 & Gaussian & 2 & 1.157e-07 & [6.848e-08, 1.806e-07] \\
low & 50 & Gaussian & 8 & 2.856e-08 & [1.705e-08, 4.046e-08] \\
low & 50 & Gaussian & 32 & 7.193e-09 & [4.243e-09, 1.027e-08] \\
low & 50 & Rademacher & 2 & 2.588e-08 & [8.015e-09, 4.475e-08] \\
low & 50 & Rademacher & 8 & 6.200e-09 & [2.034e-09, 1.148e-08] \\
low & 50 & Rademacher & 32 & 1.556e-09 & [4.686e-10, 2.983e-09] \\
long & 10 & Gaussian & 2 & 1.407e-05 & [1.231e-05, 1.676e-05] \\
long & 10 & Gaussian & 8 & 3.522e-06 & [2.959e-06, 4.388e-06] \\
long & 10 & Gaussian & 32 & 8.806e-07 & [7.080e-07, 1.069e-06] \\
long & 10 & Rademacher & 2 & 3.836e-06 & [3.107e-06, 4.673e-06] \\
long & 10 & Rademacher & 8 & 9.719e-07 & [8.110e-07, 1.213e-06] \\
long & 10 & Rademacher & 32 & 2.331e-07 & [1.977e-07, 2.797e-07] \\
long & 50 & Gaussian & 2 & 2.434e-06 & [1.748e-06, 3.182e-06] \\
long & 50 & Gaussian & 8 & 6.223e-07 & [4.397e-07, 7.653e-07] \\
long & 50 & Gaussian & 32 & 1.547e-07 & [1.070e-07, 1.950e-07] \\
long & 50 & Rademacher & 2 & 1.075e-06 & [6.689e-07, 1.605e-06] \\
long & 50 & Rademacher & 8 & 2.716e-07 & [1.780e-07, 3.970e-07] \\
long & 50 & Rademacher & 32 & 6.835e-08 & [4.474e-08, 1.046e-07] \\
stress & 10 & Gaussian & 2 & 3.350e-05 & [3.051e-05, 3.649e-05] \\
stress & 10 & Gaussian & 8 & 8.262e-06 & [7.296e-06, 9.575e-06] \\
stress & 10 & Gaussian & 32 & 2.050e-06 & [1.827e-06, 2.352e-06] \\
stress & 10 & Rademacher & 2 & 4.510e-06 & [4.023e-06, 5.437e-06] \\
stress & 10 & Rademacher & 8 & 1.141e-06 & [9.669e-07, 1.351e-06] \\
stress & 10 & Rademacher & 32 & 2.806e-07 & [2.440e-07, 3.421e-07] \\
stress & 50 & Gaussian & 2 & 7.471e-06 & [6.926e-06, 7.885e-06] \\
stress & 50 & Gaussian & 8 & 1.887e-06 & [1.790e-06, 2.059e-06] \\
stress & 50 & Gaussian & 32 & 4.733e-07 & [4.469e-07, 5.197e-07] \\
stress & 50 & Rademacher & 2 & 7.426e-07 & [5.281e-07, 9.103e-07] \\
stress & 50 & Rademacher & 8 & 1.868e-07 & [1.362e-07, 2.293e-07] \\
stress & 50 & Rademacher & 32 & 4.708e-08 & [3.412e-08, 5.920e-08] \\
\bottomrule
\end{longtable}
\par\medskip\noindent\footnotesize Each row retains all eight streams. Within each saved trial, the unbiased variance across 64 independent probe banks is averaged over 64 fixed broad-audit states; the table then reports the mean and observed range across the eight trial statistics. The traced quantity is $\operatorname{tr}(\Sigma\Sigma^\top D_x^2V)$, before the generator factor $1/2$. Gaussian uses covariance-distributed directions. Rademacher uses idiosyncratic directions and evaluates the common-shock contraction exactly. Probe counts 2, 8 and 32 use nested prefixes within each bank. The fitted HJB models use exact diffusion contractions; neither these probe outcomes nor the untouched audit selected a checkpoint. The numbers describe the frozen finite banks and introduce no additional probability event.
\endgroup


\subsection{Economic gains and complete process work}
Table~\ref{tab:r16-robustness-methods} gives all forty mean method gains and
their full-adapted-class regret upper bounds. Every mean lower endpoint
exceeds .0005. The complete direct comparison is
Table~\ref{tab:r16-robustness-comparisons} in the main article. Only the two
low-volatility, fifty-dimensional NBO--HJB comparisons have lower endpoints
above $10^{-4}$; the other thirty direct intervals cross zero. The intervals
are conditional on the complete declared eight-stream policy populations.
They do not assert a law for unseen optimizer randomness or pool across
economic designs.

\begingroup\small
\begin{longtable}{lrrrrr}
\caption{Prospective robustness over economic calibrations}\label{tab:r16-robustness-methods}\\
\toprule
Design & $d$ & Method & Mean gain & Simultaneous interval & Regret bound \\
\midrule
\endfirsthead
\toprule
Design & $d$ & Method & Mean gain & Simultaneous interval & Regret bound \\
\midrule
\endhead
high & 10 & NBO & 0.001270 & [0.000908, 0.001632] & 0.026006 \\
high & 10 & Raw & 0.001270 & [0.000908, 0.001632] & 0.026006 \\
high & 10 & Direct policy & 0.001265 & [0.000903, 0.001627] & 0.026011 \\
high & 10 & HJB greedy & 0.001262 & [0.000900, 0.001624] & 0.026014 \\
high & 10 & HJB actor & 0.001259 & [0.000898, 0.001621] & 0.026016 \\
high & 50 & NBO & 0.001347 & [0.000927, 0.001768] & 0.036873 \\
high & 50 & Raw & 0.001349 & [0.000929, 0.001770] & 0.036871 \\
high & 50 & Direct policy & 0.001349 & [0.000928, 0.001769] & 0.036872 \\
high & 50 & HJB greedy & 0.001368 & [0.000947, 0.001789] & 0.036853 \\
high & 50 & HJB actor & 0.001308 & [0.000887, 0.001729] & 0.036913 \\
low & 10 & NBO & 0.001463 & [0.001196, 0.001731] & 0.020406 \\
low & 10 & Raw & 0.001463 & [0.001196, 0.001731] & 0.020406 \\
low & 10 & Direct policy & 0.001462 & [0.001194, 0.001729] & 0.020408 \\
low & 10 & HJB greedy & 0.001440 & [0.001172, 0.001707] & 0.020430 \\
low & 10 & HJB actor & 0.001440 & [0.001172, 0.001708] & 0.020430 \\
low & 50 & NBO & 0.001604 & [0.001289, 0.001918] & 0.029761 \\
low & 50 & Raw & 0.001604 & [0.001289, 0.001918] & 0.029761 \\
low & 50 & Direct policy & 0.001601 & [0.001287, 0.001916] & 0.029763 \\
low & 50 & HJB greedy & 0.001088 & [0.000771, 0.001406] & 0.030280 \\
low & 50 & HJB actor & 0.001082 & [0.000764, 0.001400] & 0.030286 \\
long & 10 & NBO & 0.008497 & [0.006871, 0.010123] & 0.191272 \\
long & 10 & Raw & 0.008495 & [0.006869, 0.010121] & 0.191274 \\
long & 10 & Direct policy & 0.008477 & [0.006851, 0.010103] & 0.191292 \\
long & 10 & HJB greedy & 0.008437 & [0.006811, 0.010064] & 0.191332 \\
long & 10 & HJB actor & 0.008425 & [0.006799, 0.010052] & 0.191344 \\
long & 50 & NBO & 0.008272 & [0.006298, 0.010245] & 0.293097 \\
long & 50 & Raw & 0.008275 & [0.006301, 0.010248] & 0.293094 \\
long & 50 & Direct policy & 0.008266 & [0.006292, 0.010239] & 0.293103 \\
long & 50 & HJB greedy & 0.008008 & [0.006032, 0.009985] & 0.293363 \\
long & 50 & HJB actor & 0.007948 & [0.005972, 0.009924] & 0.293423 \\
stress & 10 & NBO & 0.005733 & [0.004616, 0.006850] & 0.176182 \\
stress & 10 & Raw & 0.005732 & [0.004615, 0.006849] & 0.176183 \\
stress & 10 & Direct policy & 0.005726 & [0.004609, 0.006842] & 0.176188 \\
stress & 10 & HJB greedy & 0.005712 & [0.004595, 0.006829] & 0.176202 \\
stress & 10 & HJB actor & 0.005695 & [0.004578, 0.006811] & 0.176219 \\
stress & 50 & NBO & 0.007444 & [0.006119, 0.008770] & 0.247977 \\
stress & 50 & Raw & 0.007448 & [0.006122, 0.008773] & 0.247974 \\
stress & 50 & Direct policy & 0.007425 & [0.006100, 0.008750] & 0.247996 \\
stress & 50 & HJB greedy & 0.007468 & [0.006143, 0.008794] & 0.247953 \\
stress & 50 & HJB actor & 0.007242 & [0.005917, 0.008568] & 0.248179 \\
\bottomrule
\end{longtable}
\par\medskip\noindent\footnotesize All four economic designs and all eight pre-drawn streams are retained. Each method receives 8192 independent confirmation paths on 2048 cells. Intervals share a prospectively fixed familywise allocation $\alpha=0.01$ over 648 two-sided events. The full-adapted-class regret bound is displayed beside each gain; positive improvement does not imply near optimality. Endpoints are rounded outward.
\endgroup


Table~\ref{tab:r16-robustness-work} reports the complete process work for
all outputs and the number of individual-policy final lower endpoints
exceeding .0005. A pooled positive mean is distinct from an individual
policy's target certificate. In the high-volatility, ten-dimensional cell,
all five pooled lower endpoints exceed .0005 while no individual policy's
protected final lower endpoint does. In dimension fifty, the corresponding
individual counts are 2 of 8 for NBO, 4 for Raw, 4 for direct policy, 8 for
HJB's greedy deployment, and zero for its distilled actor. All of these
outcomes remain in the comparison; a policy is not dropped because its
final certificate does not attain the target.

The recorded costs include every fixed-envelope fit, independent selection,
action construction, confirmation, and outer process overhead. Both HJB
outputs retain their shared prerequisites and their own deployment cost.
Their common construction is not divided to manufacture a cheaper actor.
These are environment-specific observed times, with complete operation
counters available in the report. They are not equal-FLOP comparisons,
runtime confidence intervals, or newly observed online stopping times.

\begingroup\small
\begin{longtable}{lrrrrr}
\caption{Complete fixed-envelope work and final target attainment}\label{tab:r16-robustness-work}\\
\toprule
Design & $d$ & Method & Mean seconds & Final target & Fallbacks \\
\midrule
\endfirsthead
\toprule
Design & $d$ & Method & Mean seconds & Final target & Fallbacks \\
\midrule
\endhead
high & 10 & NBO & 38.98 & 0/8 & 0 \\
high & 10 & Raw & 29.94 & 0/8 & 0 \\
high & 10 & Direct policy & 163.36 & 0/8 & 0 \\
high & 10 & HJB greedy & 108.91 & 0/8 & 0 \\
high & 10 & HJB actor & 85.70 & 0/8 & 0 \\
high & 50 & NBO & 96.38 & 2/8 & 0 \\
high & 50 & Raw & 86.11 & 4/8 & 0 \\
high & 50 & Direct policy & 279.87 & 4/8 & 0 \\
high & 50 & HJB greedy & 287.72 & 8/8 & 0 \\
high & 50 & HJB actor & 227.84 & 0/8 & 0 \\
low & 10 & NBO & 40.19 & 8/8 & 0 \\
low & 10 & Raw & 30.57 & 8/8 & 0 \\
low & 10 & Direct policy & 170.31 & 8/8 & 0 \\
low & 10 & HJB greedy & 117.10 & 8/8 & 0 \\
low & 10 & HJB actor & 92.77 & 8/8 & 0 \\
low & 50 & NBO & 102.27 & 8/8 & 0 \\
low & 50 & Raw & 91.29 & 8/8 & 0 \\
low & 50 & Direct policy & 310.13 & 8/8 & 0 \\
low & 50 & HJB greedy & 297.44 & 5/8 & 0 \\
low & 50 & HJB actor & 240.86 & 5/8 & 0 \\
long & 10 & NBO & 41.85 & 8/8 & 0 \\
long & 10 & Raw & 32.38 & 8/8 & 0 \\
long & 10 & Direct policy & 179.86 & 8/8 & 0 \\
long & 10 & HJB greedy & 119.14 & 8/8 & 0 \\
long & 10 & HJB actor & 93.89 & 8/8 & 0 \\
long & 50 & NBO & 99.13 & 8/8 & 0 \\
long & 50 & Raw & 88.43 & 8/8 & 0 \\
long & 50 & Direct policy & 304.95 & 8/8 & 0 \\
long & 50 & HJB greedy & 277.63 & 8/8 & 0 \\
long & 50 & HJB actor & 225.88 & 8/8 & 0 \\
stress & 10 & NBO & 42.27 & 8/8 & 0 \\
stress & 10 & Raw & 32.58 & 8/8 & 0 \\
stress & 10 & Direct policy & 185.96 & 8/8 & 0 \\
stress & 10 & HJB greedy & 124.05 & 8/8 & 0 \\
stress & 10 & HJB actor & 97.34 & 8/8 & 0 \\
stress & 50 & NBO & 100.17 & 8/8 & 0 \\
stress & 50 & Raw & 89.80 & 8/8 & 0 \\
stress & 50 & Direct policy & 306.69 & 8/8 & 0 \\
stress & 50 & HJB greedy & 282.68 & 8/8 & 0 \\
stress & 50 & HJB actor & 226.49 & 8/8 & 0 \\
\bottomrule
\end{longtable}
\par\medskip\noindent\footnotesize All streams enter the mean, including failed training and nonattainment. Target is 0.0005 at the single final decision. HJB greedy and its distilled actor share the complete four-candidate construction; each is charged that prerequisite and its own deployment, with all launch and finalization overhead. These are environment-specific measured costs, not equal FLOPs.
\endgroup

